\documentclass[../main.tex]{subfiles}
\begin{document}

El ransomware se ha convertido en una de las amenazas más graves para
organizaciones de todo tipo. Los ataques modernos no se limitan a cifrar
los datos de la víctima: incluyen sistemáticamente la destrucción de las
copias de seguridad antes del cifrado, eliminando la principal vía de
recuperación. Frente a este escenario, la industria propone el
almacenamiento inmutable como mecanismo de protección, pero la literatura
académica carece de comparativas experimentales que cuantifiquen su impacto
real en el tiempo de recuperación.

Este Trabajo de Fin de Grado presenta el diseño, implementación y evaluación
de un laboratorio de simulación de ransomware basado íntegramente en
contenedores Docker. El entorno reproduce una infraestructura empresarial
mínima compuesta por un servidor de ficheros, una \textit{workstation}, un
servidor de \textit{backups} SFTP, un sistema de almacenamiento de objetos
con Object Lock (MinIO) y un SIEM completo (Wazuh~4.7.5 con OpenSearch).

El diseño experimental sigue un esquema factorial de cuatro escenarios de
ataque por tres estrategias de recuperación y diez repeticiones, totalizando
120~ejecuciones completamente automatizadas. Los escenarios cubren las
tácticas más relevantes de la matriz MITRE ATT\&CK: cifrado masivo de
ficheros (T1486), destrucción activa de copias de seguridad (T1490),
movimiento lateral por SSH (T1021.004) y exfiltración de datos (T1041).
Las tres estrategias evaluadas son: ausencia de copia de seguridad,
repositorio Restic sobre SFTP convencional y repositorio Restic sobre MinIO
con Object Lock COMPLIANCE activo durante 30~días.

Los resultados muestran que la hipótesis principal de reducción del MTTR
en al menos un 30\,\% no se confirma estadísticamente: cuando el repositorio
está disponible, ambas estrategias restauran con velocidad equivalente. Sin
embargo, el escenario de destrucción activa revela la diferencia fundamental:
la estrategia SFTP falla en el 100\,\% de los casos al quedar el repositorio
destruido, mientras que Object Lock bloquea todos los intentos de borrado y
mantiene una tasa de éxito del 100\,\% (test exacto de Fisher $p < 0.001$).
La ventaja del almacenamiento inmutable no es de velocidad sino de
disponibilidad garantizada.

Los tiempos de detección mediante Wazuh oscilan entre 5.7 y 41.9~segundos,
frente a los más de 200~días reportados para la detección manual, confirmando
que el SIEM transforma cualitativamente la capacidad de respuesta ante
incidentes.

Como contribución adicional, se valida el laboratorio con el \textit{framework}
Atomic Red Team, ejecutando las cuatro técnicas ATT\&CK sobre el entorno real.
Todas son detectadas correctamente por las reglas desarrolladas. El laboratorio
se publica como software libre para facilitar su reproducción.

\end{document}
